\section{Compact moduli and their markings}\label{sec:geometry}

We recall precisely the geometric results about the family that will
be used. Their proofs are in \cite[\S\S4--7]{wuebben17}. All moduli
spaces in this paper refer to this fixed complex K3 surface, rather
than to an unspecified surface with the same intersection matrix.

\subsection{Central factors and compact contractions}
The embedding \eqref{eq:embedding} is primitive, and the lattice
\eqref{eq:NS} contains no roots. Indeed a root would give
$x^2-3y^2-4yz-4z^2=-1$, which is impossible modulo four, whether $y$
is even or odd. The K3 realization and the choice of ample component
in \cite[\S4.1]{wuebben17} therefore give the stated $X,H$.
Let $L_0=\OO_X$, $L_1=\OO_X(D_1)$ and $L_2=\OO_X(D_2)$. Then
\begin{equation}\label{eq:factor-pairing}
 ((a_i,a_j))_{i,j=0}^2=
 \begin{pmatrix}-2&1&2\\1&-2&1\\2&1&-2\end{pmatrix},
 \qquad v^2=2,\quad w^2=0.
\end{equation}
At $H$ the cross Hom and second Ext groups of these distinct lines
vanish. Riemann--Roch gives cross-extension dimensions $1,2,1$.

The complete central factor and charge calculation is
\cite[Lemma~4.1]{wuebben17}. For vector $v$ its polystable objects
are the stable bundles, $L_1\oplus F$ with $F$ stable of vector $w$,
and $L_0\oplus L_1\oplus L_2$. For vector $w$ the only strictly
polystable object is $L_0\oplus L_2$. Every central semistable sheaf
in both problems is locally free. The same calculation excludes all
positive lost-charge ideal objects. Surface gauge compactness then
excludes such objects throughout a sufficiently small neighborhood
of $H$, including its walls \cite[Lemma~3.2 and \S4.2]{wuebben17}.
Thus the generic chamber spaces are full compact Gieseker moduli,
consisting entirely of slope-stable bundles, and have deformation
type K3$^{[2]}$.

Here the assertion about positive charge is stronger than the absence
of singular sheaves in a particular chamber. It gives smooth gauge
subsequences when polarizations approach a wall. This is used below
to identify whole exceptional fibres with finite quotient fibres.

The auxiliary generic $w$-moduli is a smooth compact K3 surface $S$.
Its central map $S\to T$ contracts one $(-2)$-curve $C$ to the
class of $L_0\oplus L_2$. Write $S^\circ=S\setminus C$.
The normal projective central target $Z$ has the disjoint strata
\begin{equation}\label{eq:strata}
 Z=U\sqcup S^\circ\sqcup\{o\},
\end{equation}
where $U$ is the connected dense stable open set. Its reduced singular
surface is $T$. Each projective map $\pi_\sigma:M_\sigma\to Z$
is an isomorphism over $U$, has a projective-line fibre over every
point of $S^\circ$, and has the exceptional configuration described
below over $o$. Projective and fixed-$H$ gauge topologies agree.

We use the full convergent classifying charts of
\cite[\S\S2,4]{wuebben17}. At $o$ the effective compact group is
$K=\mathrm U(1)^3/\mathrm U(1)$. With single arrows
$z:1\to0$, $w_1:1\to2$ and a double arrow $u:0\to2$, and
reverse arrows $z',w_1',\varphi$, the equations are
\begin{equation}\label{eq:quiver}
 \begin{gathered}
 zz'-\varphi(u)=0,\qquad w_1w_1'+\varphi(u)=0,\\
 \mu_0=|z|^2-|z'|^2-|u|^2+|\varphi|^2,\qquad
 \mu_2=|w_1|^2-|w_1'|^2+|u|^2-|\varphi|^2,\\
 \mu_1=-\mu_0-\mu_2,\qquad \mu=\zeta.
 \end{gathered}
\end{equation}
The norms are the central harmonic norms after choosing orthonormal
coordinates. These charts commute with the projective-to-affine
quotient and the central graded-class map. They cover whole fibres
for small parameters, including moving fibres whose bases tend to
$o$. Along $S^\circ$ the transverse quotient is the two-arrow
$A_1$ model, together with its smooth surface tangent factor.
The finite wall list is exactly $\zeta_i=0$.

The whole-fibre statement has two parts. No-charge compactness and
normalized limiting inclusion and quotient maps force a sequence
over a given graded class to converge to that class. On the finite
side, vanishing cycle invariants give an acyclic arrow support.
Summing moments over its source and sink subsets bounds every arrow
norm squared by a fixed multiple of $|\zeta|$. Thus its complete
exceptional core lies in the same convergent chart. These are the
coverage results proved in \cite[\S4.4]{wuebben17}; a formal germ
alone would not have this consequence.

\subsection{Cycle classes and integral witnesses}
Let $\theta_v:v^\perp\to H^2(M_\sigma,\Z)$ be the Mukai
homomorphism, in the convention
\begin{equation}\label{eq:theta}
 \theta_v(z)=
 \left[p_{M*}\bigl(\ch(\cE)\sqrt{\td_X}\,z^\vee\bigr)\right]_2.
\end{equation}
Here $p_M:X\times M_\sigma\to M_\sigma$, $\cE$ is a universal
class and $(r,c,s)^\vee=(r,-c,s)$. The map is independent of twisting
$\cE$ by a line bundle from $M_\sigma$. We identify curve classes
with rational elements of $v^\perp$ through the Mukai pairing and
$\theta_v$. In chamber $A$ the primitive line classes in the two
planes are
\begin{equation}\label{eq:curves}
 c_0=v/2-a_0,\qquad c_2=v/2-a_2,\qquad
 f=c_0+c_2=a_1.
\end{equation}
The last class is the proper reduced specialization of a moving
$\PP^1$ fibre. The planes meet transversely once. Their normal bundles
are their cotangent bundles, and their classes satisfy
\begin{equation}\label{eq:plane-formula}
 [P_i]=\tfrac12\theta_v(c_i)^2+\tfrac1{24}c_2(M_\sigma),\qquad
 \left(\int [P_i][P_j]\right)_{i,j=0,2}
 =\begin{pmatrix}3&1\\1&3\end{pmatrix}.
\end{equation}
These are the geometric identifications of
\cite[\S\S5--6]{wuebben17}; the plane formula uses the classical
K3$^{[2]}$ intersection theory of \cite{ht}.

We will need integral, rather than rational, independence. Define
\begin{equation}\label{eq:witness}
 x=(0,f_2-f_3,0),\qquad y_0=(0,-3f_2,-1)\quad\text{in }v^\perp.
\end{equation}
The universal extensions on the planes give, with their positive
hyperplane classes $h_0,h_2$,
\begin{equation}\label{eq:witness-restrictions}
 \begin{aligned}
 \theta_v(x)|_{P_0\cup P_2}&=(0,h_2),&
 \theta_v(y_0)|_{P_0\cup P_2}&=(-h_0,-h_2),\\
 \theta_v(x)^2|_{P_0\cup P_2}&=(0,h_2^2),&
 \theta_v(y_0)^2|_{P_0\cup P_2}&=(h_0^2,h_2^2).
 \end{aligned}
\end{equation}
Both restriction matrices are unimodular. They prove that the pairs
of plane and line classes embed primitively, and that the two
evaluation maps onto $\Z^2$ are surjective.

\subsection{Universal indices and twisting}
The $v$-moduli spaces admit ordinary universal families, unique up
to a base line bundle, since $\chi(v)=-1$. The auxiliary $w$-moduli
has a nonzero algebraic Brauer obstruction of order two. We retain
local twisted universal families there. This obstruction is not
removed by constancy of the fibre cohomology of its projective
extension family.

Identify $K^0_{\mathrm{top}}(X)$ with the Mukai lattice by its Mukai
vector. This convention lets a symbol such as $y=(0,-f_2,0)$ denote
both an integral vector and the corresponding topological class;
$\theta_v$ takes the vector, whereas a family index takes the class.
For any $z\in K^0_{\mathrm{top}}(X)$ put
\begin{equation}\label{eq:normalized-index}
 \begin{gathered}
 I_z(\cE)=p_{M!}(p_X^*z^\vee\cdot\cE),\qquad y=(0,-f_2,0),\\
 \cE^y=\cE\otimes p_M^*(\det I_y(\cE))^{-1},\qquad
 I_z^y=I_z(\cE^y).
 \end{gathered}
\end{equation}
The exclamation mark denotes the complex-oriented topological
$K$-theory pushforward. Since $\chi(v,y)=-(v,y)=1$, the $y$-index
has rank one. The normalization cancels any base line twist of
$\cE$, and $\det I_y^y$ is trivial. Its restrictions to the planes
have the particularly useful form
\begin{equation}\label{eq:index-plane}
 c(I_z^y)|_{P_i}=(1+h_i)^{m_i(z)},\qquad
 m_i(z)=\chi(L_i,z)=-(a_i,z),\quad i=0,2.
\end{equation}
This is the normalized extension calculation in
\cite[\S6.4]{wuebben17}; it also holds for virtual and nonalgebraic
topological classes. Section~\ref{sec:ring} computes all these
indices globally from the universal subscheme on $X^{[2]}$.
